\documentclass[10pt,letterpaper]{article}
\usepackage{cornell-notes}

\title{Clause 29.11.02.04: Remote Time Zone Database Support}
\author{}
\date{}
\setDocTitle{Clause 29.11.02.04: Remote Time Zone Database Support}
\setDocSubtitle{C++ 2024}
\setDocOwner{C++ 2024 Cornell Notes}
\setDocDate{\today}
\setCornellCollection{C++ 2024 Cornell Notes}
\setCornellUnitType{Clause}
\setCornellUnitNumber{29.11.02.04}
\setCornellUnitTitle{Remote Time Zone Database Support}
\hypersetup{pdftitle={Clause 29.11.02.04: Remote Time Zone Database Support},pdfsubject={Cornell notes},pdfauthor={}}

\begin{document}
\maketitle
\begin{CornellOverviewBox}{Overview}
\textbf{Classification:} Standard library - Normative\\[2pt]
\textbf{Learning objective:} Understand the rules, terminology, and semantic consequences associated with Remote time zone database support.
\end{CornellOverviewBox}\begin{CornellSummaryBox}{Summary}
{\sffamily\bfseries\color{Accent} Summary}\par\vspace{3pt}Section 29.11.2.4 focuses on Remote time zone database support. Apply the specified interface together with its mandates, effects, result conditions, exception behavior, and complexity constraints. When applying it, preserve the distinction between mandatory requirements, recommendations, permissions, and behavior deliberately left undefined, unspecified, or implementation-defined.
\end{CornellSummaryBox}

\end{document}
